\documentclass[../../main.tex]{subfiles}

\begin{document}

\section{凸泛函的最小化}\label{section:ch08-凸泛函最小化}

\begin{definition}[凸集与闭集]\label{definition:ch08-凸集闭集}
赋范线性空间 $X$ 中的子集 $C$ 称为凸的, 若对所有 $f$, $g\in C$ 和 $\lambda\in[0,1]$,
\begin{align*}
\lambda f+(1-\lambda)g\in C.
\end{align*}
$X$ 中的子集 $C$ 称为闭的, 若 $C$ 中每个强收敛序列的极限属于 $C$.
\end{definition}

\begin{definition}[连续与凸泛函]\label{definition:ch08-连续凸泛函}
定义在赋范线性空间 $X$ 的子集 $C$ 上的实值泛函 $\Phi$ 称为连续的, 若对 $C$ 中每个在 $X$ 中强收敛于 $f\in C$ 的序列 $\{f_{n}\}$,
\begin{align*}
\lim_{n\to\infty}\Phi(f_{n})=\Phi(f).
\end{align*}
泛函 $\Phi$ 称为凸的, 若对所有 $f$, $g\in C$ 和 $\lambda\in[0,1]$,
\begin{align*}
\Phi(\lambda f+(1-\lambda)g)\leqslant\lambda\Phi(f)+(1-\lambda)\Phi(g).
\end{align*}
\end{definition}

\begin{lemma}\label{lemma:ch08-弱收敛极限属于凸闭集}
令 $1<p<\infty$ 而 $C$ 是 $L^{p}(E)$ 中的有界闭凸集. 若在 $L^{p}(E)$ 中 $\{f_{n}\rightharpoonup f\}$ 且每个 $f_{n}\in C$, 则 $f\in C$.
\end{lemma}

\begin{proof}
根据 Helly 定理, $C$ 是弱序列紧的. 由于 $C$ 是凸的闭集, 它在弱收敛下保持闭性 (这是 Mazur 引理的结论: 凸闭集是弱闭的). 因此 $f\in C$.
\end{proof}

\begin{theorem}[凸泛函最小化定理]\label{theorem:ch08-凸泛函最小化}
令 $1<p<\infty$, $C$ 是 $L^{p}(E)$ 中的非空有界闭凸集, 而 $\Phi$ 是 $C$ 上的连续凸泛函. 则 $\Phi$ 在 $C$ 上取到最小值, 即存在 $f_{0}\in C$ 使得
\begin{align*}
\Phi(f_{0})=\min_{f\in C}\Phi(f).
\end{align*}
\end{theorem}

\begin{proof}
令 $\{f_{n}\}$ 为极小化序列, 即 $\lim_{n\to\infty}\Phi(f_{n})=\inf_{f\in C}\Phi(f)$. 由于 $C$ 有界, 由 Helly 定理, 存在子序列 $\{f_{n_{k}}\}$ 弱收敛于某个 $f_{0}\in L^{p}(E)$. 由引理 16, $f_{0}\in C$. 由凸性, $\Phi$ 是弱下半连续的: 若在 $L^{p}(E)$ 中 $\{f_{n_{k}}\rightharpoonup f_{0}\}$, 则
\begin{align*}
\Phi(f_{0})\leqslant\liminf_{k\to\infty}\Phi(f_{n_{k}})=\inf_{f\in C}\Phi(f).
\end{align*}
因此 $\Phi(f_{0})=\inf_{f\in C}$, 即 $f_{0}$ 是最小值点.
\end{proof}

\begin{corollary}\label{corollary:ch08-范数最小化}
令 $1<p<\infty$ 而 $C$ 是 $L^{p}(E)$ 中的非空有界闭凸集. 则存在 $f_{0}\in C$ 使得
\begin{align*}
\|f_{0}\|_{p}=\min_{f\in C}\|f\|_{p}.
\end{align*}
\end{corollary}

\begin{proof}
范数 $\Phi(f)=\|f\|_{p}$ 是连续且凸的.
\end{proof}

\begin{exercise}\label{exercise:ch08-38}
证明: 若 $C$ 是凸集, 则 $C$ 中序列的凸组合也属于 $C$.
\end{exercise}

\begin{exercise}\label{exercise:ch08-39}
证明: 若 $\Phi$ 是凸泛函, 则 $\Phi$ 是弱下半连续的.
\end{exercise}

\begin{exercise}\label{exercise:ch08-40}
证明: 若 $C$ 是闭凸集, 则 $C$ 是弱闭的 (Mazur 引理).
\end{exercise}

\begin{exercise}\label{exercise:ch08-41}
令 $1<p<\infty$ 而 $C$ 是 $L^{p}(E)$ 中的非空闭凸集 (不一定有界). 证明: 若 $\lim_{\|f\|_{p}\to\infty}\Phi(f)=\infty$, 则 $\Phi$ 在 $C$ 上取到最小值.
\end{exercise}

\begin{exercise}\label{exercise:ch08-42}
举例说明: 若 $C$ 不是凸的, 则有界闭集上的连续泛函不一定取到最小值.
\end{exercise}

\begin{exercise}\label{exercise:ch08-43}
令 $1<p<\infty$. 证明: 若 $C$ 是 $L^{p}(E)$ 中的非空闭凸集, 则对每个 $f\in L^{p}(E)$, 存在唯一的 $f_{0}\in C$ 使得 $\|f-f_{0}\|_{p}=\min_{g\in C}\|f-g\|_{p}$.
\end{exercise}

\begin{exercise}\label{exercise:ch08-44}
(Banach-Saks 定理) 令 $1<p<\infty$ 而 $\{f_{n}\}$ 是 $L^{p}(E)$ 中的有界序列. 证明存在子序列 $\{f_{n_{k}}\}$ 和函数 $f\in L^{p}(E)$, 使得算术平均序列
\begin{align*}
g_{k}=\frac{1}{k}\sum_{j=1}^{k}f_{n_{j}}
\end{align*}
在 $L^{p}(E)$ 中强收敛于 $f$.
\end{exercise}

\begin{exercise}\label{exercise:ch08-45}
用 Banach-Saks 定理证明凸泛函的弱下半连续性.
\end{exercise}

\end{document}